\section{Introduction}

\subsection{Background}
Recent advances in Large Language Models (LLMs) have enabled a new generation of autonomous agents capable of reasoning, planning, and executing tasks with minimal human intervention. As these agents become more capable, the research community and industry have moved from single-agent systems to multi-agent collaborative frameworks, where multiple specialized agents work together to solve complex tasks.

Multi-agent collaboration typically follows one of two patterns. In the first pattern, a \textbf{monolithic agent} handles all tasks internally using tools and memory. In the second pattern, a \textbf{multi-agent system} decomposes a complex task into subtasks, assigning each to a specialized agent that focuses on a narrow domain.

However, a critical architectural question arises: how should multiple agent runtimes be organized and coordinated? Current frameworks, including AutoGen~\citep{wu2023autogen}, LangChain Agents~\citep{langchain}, CrewAI~\citep{crewai}, MetaGPT~\citep{hong2024metagpt}, and ChatDev~\citep{qian2024chatdev}, predominantly adopt a \textbf{Single-Runtime Multi-Agent} architecture, where all agents share the same process space, memory, and tool registry.

\subsection{Problem Statement}
Consider a practical scenario: a researcher needs OpenClaw for daily task management, Claude Code for code writing, and a domain-specific agent (drsai) for physics analysis. The user must coordinate these specialized agents individually---which agent should handle this task? How should results be aggregated? This coordination burden grows exponentially as the number of agents increases.

Current multi-agent frameworks, while theoretically supporting multiple agents, fail to address this fundamental usability problem:

\textbf{The Coordination Problem}: Users must manually decompose tasks and assign them to specialized agents. There is no native mechanism for an Orchestrator to act as the single interface between the user and their agent team.

\textbf{Isolation Problem}: Even if we could coordinate agents, Single-Runtime architectures place all agents in the same process space. A bug or crash in one agent cascades to others.

\textbf{Scalability Problem}: Adding a new agent type requires modifying the main process code, as all agents must be registered in the same tool registry.

\textbf{Security Problem}: In March 2026, CNCERT and the China Cyberspace Security Association issued a joint security advisory on OpenClaw~\citep{cncert2026openclaw}. The advisory identified four major security risks and explicitly recommended \textbf{container-based isolation} as a primary mitigation strategy.

\textbf{Interoperability Problem}: Different agent runtimes (e.g., OpenClaw, Claude Code) have no native mechanism to coordinate, as each defines its own tool and skill interfaces.

\subsection{Contributions}
This paper makes the following contributions:

\begin{enumerate}
    \item \textbf{Agent-as-Container Abstraction}: We propose treating each agent runtime as an isolated container managed by Kubernetes, enabling true runtime isolation, independent scaling, and fast lifecycle management.
    
    \item \textbf{Multi-Runtime Multi-Agent Architecture}: We design and implement a hierarchical architecture where an Orchestrator agent coordinates multiple specialized agent runtimes, each running in its own container.
    
    \item \textbf{Skill as the Standardized Coordination Interface}: We introduce Skill as a first-class abstraction that allows an Orchestrator to invoke any agent runtime without knowledge of its internal APIs. This enables the Orchestrator to serve as the \textbf{single interface} between the user and their agent team.
    
    \item \textbf{DORA: A Working Demonstration}: We implement DORA, an Orchestrator that coordinates OpenClaw, Claude Code, and drsai through Skills. The user interacts only with DORA; DORA handles task decomposition, agent dispatch, and result aggregation.
    
    \item \textbf{Production-Grade Implementation}: We implement and deploy a containerized agent platform serving researchers, with production-ready Skills for lifecycle management and orchestration.
    
    \item \textbf{Security by Design}: The containerized architecture directly addresses the security concerns raised by CNCERT/CERT, providing hardware-level isolation between agent runtimes.
\end{enumerate}

\subsection{Paper Organization}
The rest of this paper is organized as follows. Section~\ref{sec:background} provides background and discusses related work. Section~\ref{sec:design} presents the system design. Section~\ref{sec:implementation} describes the implementation. Section~\ref{sec:evaluation} evaluates the system through experiments. Section~\ref{sec:conclusion} concludes with future directions.
